\section{Construcción de un Coding Agent desde cero}

\subsection{¿Qué es un Coding Agent?}

Un Coding Agent es un sistema de IA capaz de ejecutar tareas de programación de forma autónoma. Sus componentes centrales incluyen:

\begin{itemize}
    \item \textbf{System Prompt}: define el comportamiento y las instrucciones del Agent
    \item \textbf{User Prompt}: la solicitud personalizada del usuario
    \item \textbf{Assistant Prompt}: la respuesta del LLM
    \item \textbf{Tool Set}: el conjunto de herramientas que el Agent puede invocar
\end{itemize}

\begin{importantbox}{El ciclo central de un Coding Agent}
El flujo de trabajo básico de un Coding Agent es extremadamente simple:
\begin{enumerate}
    \item Lee la entrada del usuario desde la terminal y la agrega continuamente al historial de la conversación
    \item Informa al LLM las herramientas disponibles (Read\_file, List\_dir, Edit\_file, etc.)
    \item El LLM solicita invocar una herramienta en el momento apropiado
    \item Ejecuta la llamada a la herramienta fuera de línea y devuelve el resultado
    \item Repite este proceso hasta que la tarea se complete
\end{enumerate}
\end{importantbox}

\subsection{«Arma secreta»: la filosofía de diseño de Claude Code}

Un análisis profundo del diseño subyacente de Claude Code revela varias decisiones de ingeniería clave:

\begin{enumerate}
    \item \textbf{Inyección anticipada de context}: utiliza un prompt pequeño y preciso para realizar la precarga de context
    \item \textbf{System reminders omnipresentes}: se insertan etiquetas \texttt{<system-reminder>} en el system prompt, el user prompt, las tool calls y los tool results, para evitar que el modelo se desvíe del objetivo en conversaciones largas
    \item \textbf{Command prefix extraction}: extrae el prefijo del comando a partir de la entrada del usuario para habilitar atajos de operación
    \item \textbf{Generación de subagents}: genera subagents que procesan subtareas específicas, evitando que un único context se sobrecargue
\end{enumerate}

\begin{warningbox}{El problema del Context Drift}
En las conversaciones largas, el LLM tiende a «olvidar» las instrucciones y restricciones iniciales: a esto se le llama context drift. Claude Code contrarresta este problema insertando repetidamente system reminders en cada etapa de la conversación. Esta es una técnica clave validada mediante ingeniería.
\end{warningbox}

\subsection{Resumen de la sección}

La barrera técnica para construir un Coding Agent básico es en realidad muy baja: el núcleo es solo el ciclo de LLM + tool calling. Pero construir un Coding Agent \textbf{confiable} requiere una gran cantidad de detalles de ingeniería: gestión de context, protección contra drift, descomposición de subtareas, entre otros.

